% !TEX root = ../main.tex
\documentclass[../main.tex]{subfiles}

\begin{document}

\chapter{Classical Mechanics}
\label{chap:mechanics}

\section{Newton's Laws}

\begin{law}[Newton's second law]
  The net force acting on a particle equals the time derivative of its
  momentum:
  \begin{equation}
    \vect{F}=\dv{\vect{p}}{t}.
    \label{eq:newton-second-law}
  \end{equation}
\end{law}

For constant mass, \cref{eq:newton-second-law} becomes
\begin{equation}
  \vect{F}=m\vect{a}.
\end{equation}

\section{Dimensional Analysis}

The SI unit of force is
$\si{\newton}=\si{\kilogram\metre\per\second\squared}$.

\begin{warningbox}
  State the reference frame before writing the equations of motion. The
  standard form of Newton's second law applies directly only in an inertial
  frame.
\end{warningbox}

\end{document}
